{\Large\textbf{Disk in gas}}

Consider a thin flat disk of mass $M$ and face area $S$ at temperature $T_1$
resting initially in weightlessness in a gas of mass density $\rho$ at
temperature $T_0$ ($T_1 = 1000\ T_0$). One of the faces of the disk is covered
with a thermally insulating layer, the other face has a very good thermal
contact with the surrounding gas: gas molecules of mass $m$ obtain the
temperature of the disk during a single collision with the surface. Estimate
the initial acceleration $a_0$ and maximal speed $v_{max}$ of the disk during
its subsequent motion. Assume the heat capacity of the disk to be on the of
order $Nk_B$, where $N$ is the number of atoms in it, and $k_B$ is the
Boltzmann constant, and molar masses of the gas and the disk's material to be
of the same order. The mean free path length of molecules is much larger than
the size of the disk. Neglect any edge effects occuring at the edge of the
disk.
